\documentclass[a4paper,12pt]{article}

\usepackage[T2A]{fontenc}
\usepackage[utf8]{inputenc}
\usepackage[russian]{babel}
\usepackage{amsmath,amssymb}
\usepackage[a4paper,margin=2cm]{geometry}
\usepackage{microtype}
\usepackage{array,tabularx,booktabs}
\usepackage{enumitem}
\usepackage{xcolor}
\usepackage{hyperref}
\usepackage{tikz}

\hypersetup{
  colorlinks=true,
  linkcolor=blue!55!black,
  urlcolor=blue!55!black
}

\setlength{\parindent}{0pt}
\setlength{\parskip}{0.35em}
\setlist[enumerate]{leftmargin=*,itemsep=0.15em,topsep=0.25em}
\renewcommand{\arraystretch}{1.18}

\definecolor{accent}{RGB}{37,78,112}
\definecolor{resultcolor}{RGB}{25,107,91}

\newcommand{\errmark}{\textcolor{red!75!black}{\textbf{Ошибка:}}}
\newcommand{\ideamark}{\textcolor{accent}{\textbf{Ключевая идея:}}}
\newcommand{\resultmark}{\textcolor{resultcolor}{\textbf{Итог:}}}

\newcolumntype{Y}{>{\raggedright\arraybackslash}X}

\begin{document}

\begin{titlepage}
  \centering
  \vspace*{0.16\textheight}

  {\LARGE\bfseries Работа над ошибками по домашней работе\par}
  \vspace{1.6em}
  {\large Математика\par}
  \vspace{0.45em}
  {\large Тема: производная и свойства функции по графику\par}

  \vspace{2.4em}
  \begin{tabular}{@{}rl@{}}
    \textbf{Ученица:} & Николь Саркисьянц \\
    \textbf{Преподаватель:} & Лёвин Артём Александрович \\
    \textbf{Дата:} & 25 сентября 2026 года
  \end{tabular}

  \vfill
  {\small Лёвин Артём Александрович эксклюзивно для Николь Саркисьянц\par}
  \vspace{0.8em}
  \begin{tikzpicture}[x=1cm,y=1cm]
    \draw[accent!55, line width=0.7pt]
      (0,0) .. controls (2,0.28) and (4,-0.28) .. (6,0)
            .. controls (8,0.28) and (10,-0.28) .. (12,0);
  \end{tikzpicture}
  \vspace*{0.7cm}
\end{titlepage}

\section*{Краткая сводка}
\small
\begin{tabularx}{\linewidth}{@{}>{\raggedright\arraybackslash}p{1.4cm} Y Y >{\raggedright\arraybackslash}p{2.0cm} >{\raggedright\arraybackslash}p{2.2cm}@{}}
\toprule
\textbf{№} & \textbf{Тема} & \textbf{Суть ошибки} & \textbf{Ваш ответ} & \textbf{Верный ответ} \\
\midrule
\hyperref[task:D1.5]{Д1.5}
& Знак производной и убывание функции
& После верного перехода к участкам убывания посчитано число таких участков, а требовалось число целых значений аргумента на них.
& $4$
& $5$ \\
\addlinespace
\hyperref[task:D1.12]{Д1.12}
& Минимум функции по графику производной
& Ход решения не доведён до критерия минимума: не отмечена смена знака производной с минуса на плюс.
& не указан
& $1$ \\
\bottomrule
\end{tabularx}
\normalsize

\clearpage
\phantomsection\label{task:D1.5}
\section*{Задание Д1.5}

\textbf{Текст задания.}
На рисунке изображён график функции $y=f(x)$, определённой на интервале $(-1;13)$. Определите количество целых чисел $x$, для которых $f'(x)<0$.

\textbf{Ваш ход решения.}
Вы верно отметили условие $f'(x)<0$ и связали его с убыванием функции. После записи о том, что функция убывает, сразу указан результат $4$.

\textbf{Ваш ответ.} $4$.

\errmark{} посчитано количество промежутков убывания, а не количество целых значений $x$, для которых производная отрицательна.

\textbf{Где именно возникает ошибка.}
Последний правильный шаг --- переход от условия $f'(x)<0$ к участкам, где график функции идёт вниз слева направо. Следующий шаг должен был состоять в подсчёте целых значений аргумента внутри этих участков. Вместо этого в ответе оказалось число самих участков убывания.

\textbf{Почему это неверно.}
Условие спрашивает не о количестве промежутков, а о количестве \emph{целых чисел} $x$. Поэтому каждый найденный промежуток убывания нужно рассмотреть отдельно и выписать все целые значения, лежащие строго внутри него.

По графику функция убывает на промежутках
\[
(0;1),\qquad (2;4),\qquad (5;7),\qquad (9;13).
\]
На концах этих промежутков производная не является отрицательной: в точках локальных экстремумов она равна нулю, а число $13$ не принадлежит области определения. Поэтому границы не включаются.

\ideamark{} сначала находите все участки, где функция убывает, а затем отдельно считаете целые значения $x$ внутри этих открытых промежутков.

\textbf{Исправление от последнего правильного шага.}
Выпишем целые числа по каждому промежутку:
\[
(0;1)\colon \text{целых чисел нет},
\]
\[
(2;4)\colon 3,
\qquad
(5;7)\colon 6,
\qquad
(9;13)\colon 10,\ 11,\ 12.
\]
Итого получаем
\[
1+1+3=5.
\]

\textbf{Полное правильное решение.}
\begin{enumerate}
  \item Неравенство $f'(x)<0$ выполняется там, где функция $f$ убывает.
  \item По графику участки убывания: $(0;1)$, $(2;4)$, $(5;7)$ и $(9;13)$.
  \item Среди целых значений аргумента подходят только $3$, $6$, $10$, $11$ и $12$.
  \item Таких значений пять.
\end{enumerate}

\textbf{Самоконтроль.}
Полученный список содержит пять различных целых чисел. Каждое из них находится строго внутри участка убывания, а точки, где производная равна нулю, в список не включены.

\resultmark{} правильный ответ: $5$.

\medskip
\textbf{\textcolor{accent}{Задача на закрепление.}}
Для функции $g$ известно, что $g'(x)<0$ на промежутках $(1;4)$, $(6;9)$ и $(11;15)$, а на остальных точках области определения $(0;16)$ производная неотрицательна. Сколько целых значений $x$ из области определения удовлетворяют условию $g'(x)<0$?

\clearpage
\phantomsection\label{task:D1.12}
\section*{Задание Д1.12}

\textbf{Текст задания.}
На рисунке изображён график производной функции $f(x)$, определённой на интервале $(-8;4)$. Найдите количество точек минимума функции $f(x)$, принадлежащих отрезку $[-7;-1]$.

\textbf{Ваш ход решения.}
Вы верно определили, что нужно работать с графиком $f'(x)$ и искать точки минимума функции $f(x)$, но дальнейший анализ знака производной не записан.

\textbf{Ваш ответ.} не указан.

\errmark{} решение не завершено: не применён признак минимума по смене знака производной.

\textbf{Где именно возникает ошибка.}
Последний правильный шаг --- понимание того, что экстремумы исходной функции определяются по поведению её производной. Следующий обязательный шаг --- пройти по заданному отрезку слева направо и отметить пересечения графика $f'(x)$ с осью $Ox$, одновременно определяя направление смены знака.

\textbf{Почему это важно.}
Равенства $f'(x)=0$ самого по себе недостаточно. Минимум соответствует смене знака производной $-\!\to\!+$: слева функция убывает, справа возрастает. Смена $+\!\to\!-$, наоборот, соответствует максимуму.

На отрезке $[-7;-1]$ график производной один раз пересекает ось $Ox$ снизу вверх и один раз --- сверху вниз. Первый переход даёт минимум, второй --- максимум. Следующий переход снизу вверх находится уже правее $-1$ и не учитывается.

\ideamark{} при работе с графиком производной важен не сам факт пересечения оси $Ox$, а направление смены знака.

\textbf{Исправление от последнего правильного шага.}
На отрезке $[-7;-1]$ переход $-\!\to\!+$ встречается ровно один раз, поэтому точка минимума одна.

\textbf{Полное правильное решение.}
\begin{enumerate}
  \item Ищем на графике $f'$ переходы из отрицательных значений в положительные.
  \item На отрезке $[-7;-1]$ такой переход происходит один раз.
  \item Пересечение оси в противоположном направлении даёт максимум и не учитывается.
\end{enumerate}

\textbf{Самоконтроль.}
В найденной точке слева $f'(x)<0$, а справа $f'(x)>0$, поэтому функция сначала убывает, затем возрастает.

\resultmark{} правильный ответ: $1$.

\medskip
\textbf{\textcolor{accent}{Задача на закрепление.}}
На отрезке $[-6;5]$ производная функции $h$ меняет знак с отрицательного на положительный между $-5$ и $-4$, с положительного на отрицательный между $-2$ и $-1$, а затем снова с отрицательного на положительный между $2$ и $3$. Сколько точек минимума функции $h$ находится на этом отрезке?

\medskip
\textbf{\textcolor{accent}{Ответы к задачам на закрепление}}

\smallskip
\begin{tabular}{@{}ll@{}}
\toprule
\textbf{К заданию} & \textbf{Ответ} \\
\midrule
Д1.5 & $7$ \\
Д1.12 & $2$ \\
\bottomrule
\end{tabular}

\end{document}
